\textbf{Replacement changes the coverage target.} Exact zeros arise by censoring components below 0.025 and reclosing. The fixed transform is $T_\epsilon(y)_j=\max(y_j,\epsilon)/\sum_l\max(y_l,\epsilon)$. Left: clr methods calibrate and test on $T_\epsilon(Y)$. These same values and intervals exactly equal original-response coverage of the inverse-image region $\{y:T_\epsilon(y)\in C_\epsilon(x)\}$ by event identity; this is not another simulation. Right: those same finite clr regions are tested against original $Y$, rejecting boundary points, while the raw Euclidean reference separately calibrates and tests on original $Y$. All fits use transformed training responses. Dots are means and bars pointwise 95\% $t$ intervals across 50 independent datasets, with 1000 training, 499 calibration and 5000 iid test observations each. The same original datasets are reused across replacement values. The dashed line marks 0.90; the original-clr check has no nominal guarantee. The right-panel diagnostic is a different region from the left-panel inverse image. Replacement does not inherently prevent coverage of original zeros. No inverse-image area was estimated, and this censoring test does not establish structural-zero performance.
